\section{Introduction}

Open air traffic data have enabled a wide range of research in aviation surveillance, environmental analysis, and operational modeling. Community-driven infrastructures such as the OpenSky Network have shown that large-scale aircraft trajectory data can support reproducible research at unprecedented scale~\cite{opensky2014,strohmeier2021crowdsourced}. EUROCONTROL's Aviation Data Repository for Research (ADRR) extends this opportunity with large, structured trajectory and flight-information archives for commercial flights operating in and over Europe~\cite{eurocontrol_adrr_2025}. These data sources make it realistic to revisit trajectory forecasting from the perspective of open, large-scale, and reproducible modeling rather than narrowly curated benchmark scenarios.

However, the forecasting problem supported by such data is not identical to the problem emphasized in much of the aircraft trajectory prediction literature. Many prior works focus on short-term continuous trajectory extrapolation, four-dimensional point prediction, or terminal-area interactions~\cite{jia2022attentionlstm,ma2024attentiontcngru,li2025noiserobust,shan2025sparse}. Those formulations are valuable, but they are less well aligned with long-horizon forecasting from open operational data, where observations are noisy, route structure is highly multimodal, and exact pointwise forecasts become difficult to interpret or calibrate. For strategic traffic awareness, surveillance resilience, and high-level flow reasoning, a probabilistic regional forecast can be more appropriate than a single predicted coordinate sequence.

This paper therefore studies \emph{multi-horizon regional aircraft trajectory forecasting}. Instead of predicting one future track, we predict distributions over future H3 regions at multiple lead times. This formulation offers three advantages. First, it respects the uncertainty that naturally grows with forecast horizon. Second, it aligns with large-scale open trajectory archives whose spatial quality and operational diversity make exact long-range point regression brittle. Third, it enables calibration-aware evaluation through negative log-likelihood, Brier score, expected calibration error, and candidate-region recall under explicit area budgets.

Our starting point is an existing hierarchical H3 regional forecaster built around a GRU encoder and coarse-to-fine decoding. We keep the hierarchical spatial formulation, because it already provides a scalable way to model regional uncertainty. We then ask a narrower research question: what temporal architecture best supports multi-horizon probabilistic regional forecasting on large-scale open aviation trajectories? Based on a staged empirical study, we find that the strongest improvement comes from replacing the single summary representation with horizon-specific query decoding over Transformer encoder states, and further strengthening this with an explicit multiscale temporal branch.

The resulting model, a \emph{multiscale horizon-query Transformer}, uses one encoder over fine-scale history, one encoder over pooled coarse-scale history, and one learned query per forecast horizon. The hierarchical H3 decoder remains compatible with the original baseline, which makes the architectural changes easy to isolate through ablations. On the full dataset, the final model consistently outperforms a classical Markov baseline, the original hierarchical GRU baseline, and a single-scale Transformer horizon decoder.

The main contributions of this work are:
\begin{enumerate}
  \item We formulate large-scale multi-horizon regional aircraft trajectory forecasting on open ADRR-derived flight trajectories as a hierarchical H3 probabilistic prediction problem.
  \item We propose a multiscale horizon-query Transformer that improves regional accuracy, likelihood, Brier score, and calibration over a hierarchical GRU baseline while preserving the same coarse-to-fine output structure.
  \item We introduce a candidate-region evaluation view that jointly reports recall, candidate cell count, candidate area, and calibration rather than relying on top-k accuracy alone.
  \item We provide full-scale comparisons against a classical Markov regional transition baseline and a Transformer structural ablation, and show that lightweight H3-neighbor refinement creates a recall-versus-calibration trade-off rather than a free gain.
\end{enumerate}
